\section{Localization of triangulated categories}
\label{section-localization}

\noindent
In order to construct the derived category starting from the homotopy
category of complexes, we will use a localization process.

\begin{definition}
\label{definition-localization}
Let $\mathcal{D}$ be a pre-triangulated category. We say a multiplicative
system $S$ is {\it compatible with the triangulated structure} if
the following two conditions hold:
\begin{enumerate}
\item[MS5] For a morphism $f$ of $\mathcal{D}$ we have
$f \in S \Leftrightarrow f[1] \in S$\footnote{See Remark \ref{remark-MS5}.}.
\item[MS6] Given a solid commutative square
$$
\xymatrix{
X \ar[r] \ar[d]^s &
Y \ar[r] \ar[d]^{s'} &
Z \ar[r] \ar@{-->}[d] &
X[1] \ar[d]^{s[1]} \\
X' \ar[r] &
Y' \ar[r] &
Z' \ar[r] &
X'[1]
}
$$
whose rows are distinguished triangles with $s, s' \in S$
there exists a morphism $s'' : Z \to Z'$ in $S$ such that
$(s, s', s'')$ is a morphism of triangles.
\end{enumerate}
\end{definition}

\noindent
It turns out that these axioms are not independent of the
axioms defining multiplicative systems.

\begin{lemma}
\label{lemma-localization-conditions}
Let $\mathcal{D}$ be a pre-triangulated category.
Let $S \subset \text{Arrows}(\mathcal{D})$.
\begin{enumerate}
\item If $S$ contains all identities and MS6 holds
(Definition \ref{definition-localization}),
then every isomorphism of $\mathcal{D}$ is in $S$.
\item If MS1, MS5 (Categories, Definition
\ref{categories-definition-multiplicative-system}) and MS6 hold,
then MS2 holds.
\end{enumerate}
\end{lemma}

\begin{proof}
Assume $S$ contains all identities and MS6 holds.
Let $f : X \to Y$ be an isomorphism of $\mathcal{D}$.
Consider the diagram
$$
\xymatrix{
0 \ar[r] \ar[d]^1 &
X \ar[r]_1 \ar[d]^1 &
X \ar[r] \ar@{-->}[d] &
0[1] \ar[d]^{1[1]} \\
0 \ar[r] &
X \ar[r]^f &
Y \ar[r] &
0[1]
}
$$
The rows are distinguished triangles by Lemma \ref{lemma-third-object-zero}.
By MS6 we see that the dotted arrow exists and is in $S$, so $f$ is in $S$.

\medskip\noindent
Assume MS1, MS5, MS6.
Suppose that $f : X \to Y$ is a morphism of $\mathcal{D}$ and
$t : X \to X'$ an element of $S$. Choose a distinguished triangle
$(X, Y, Z, f, g, h)$. Next, choose a distinguished triangle
$(X', Y', Z, f', g', t[1] \circ h)$ (here we use TR1 and TR2).
By MS5, MS6 (and TR2 to rotate) we can find the dotted arrow
in the commutative diagram
$$
\xymatrix{
X \ar[r] \ar[d]^t &
Y \ar[r] \ar@{..>}[d]^{s'} &
Z \ar[r] \ar[d]^1 &
X[1] \ar[d]^{t[1]} \\
X' \ar[r] &
Y' \ar[r] &
Z \ar[r] &
X'[1]
}
$$
with moreover $s' \in S$. This proves LMS2. The proof of RMS2 is dual.
\end{proof}

\begin{remark}
\label{remark-MS5}
In the presence of MS1 and MS6, condition MS5 is equivalent to asking
$s[n] \in S$ for all $s \in S$ and $n \in \mathbf{Z}$. For example,
suppose MS5 holds, we have $s \in S$, and we want to show $s[-1] \in S$.
This isn't immediate because $s[-1][1]$ is not equal to $s$, only
isomorphic to $s$ as an arrow of $\mathcal{D}$. Still, this does imply that
$s[-1][1] = f \circ s \circ g$ for isomorphisms $f$, $g$. By
Lemma \ref{lemma-localization-conditions} (1) we find $f, g \in S$,
hence $s[-1][1] \in S$ by MS1, hence $s[-1] \in S$ by MS5.
We leave a complete proof to the reader as an exercise.
\end{remark}

\begin{lemma}
\label{lemma-triangle-functor-localize}
Let $F : \mathcal{D} \to \mathcal{D}'$ be an exact functor of
pre-triangulated categories.  Let
$$
S = \{f \in \text{Arrows}(\mathcal{D})
\mid F(f)\text{ is an isomorphism}\}
$$
Then $S$ is a saturated (see
Categories,
Definition \ref{categories-definition-saturated-multiplicative-system})
multiplicative system compatible with the
triangulated structure on $\mathcal{D}$.
\end{lemma}

\begin{proof}
We have to prove axioms MS1 -- MS6, see
Categories, Definitions \ref{categories-definition-multiplicative-system} and
\ref{categories-definition-saturated-multiplicative-system}
and
Definition \ref{definition-localization}.
MS1, MS4, and MS5 are direct from the definitions. MS6 follows from TR3 and
Lemma \ref{lemma-third-isomorphism-triangle}.
By
Lemma \ref{lemma-localization-conditions}
we conclude that MS2 holds. To finish the proof we have to show that
MS3 holds. To do this let $f, g : X \to Y$ be morphisms of $\mathcal{D}$,
and let $t : Z \to X$ be an element of $S$ such that $f \circ t = g \circ t$.
As $\mathcal{D}$ is additive this simply means that $a \circ t = 0$ with
$a = f - g$. Choose a distinguished triangle $(Z, X, Q, t, d, h)$ using TR1.
Since $a \circ t = 0$ we see by
Lemma \ref{lemma-representable-homological}
there exists a morphism $i : Q \to Y$ such that $i \circ d = a$.
Finally, using TR1 again we can choose a triangle
$(Q, Y, W, i, j, k)$. Here is a picture
$$
\xymatrix{
Z \ar[r]_t & X \ar[r]_d \ar[d]^1 & Q \ar[r] \ar[d]^i & Z[1] \\
& X \ar[r]_a & Y \ar[d]^j \\
& & W
}
$$
OK, and now we apply the functor $F$ to this diagram.
Since $t \in S$ we see that $F(Q) = 0$, see
Lemma \ref{lemma-third-object-zero}.
Hence $F(j)$ is an isomorphism by the same lemma, i.e., $j \in S$.
Finally, $j \circ a = j \circ i \circ d = 0$ as $j \circ i = 0$.
Thus $j \circ f = j \circ g$ and we see that LMS3 holds.
The proof of RMS3 is dual.
\end{proof}

\begin{lemma}
\label{lemma-homological-functor-localize}
Let $H : \mathcal{D} \to \mathcal{A}$ be a homological functor between a
pre-triangulated category and an abelian category. Let
$$
S = \{f \in \text{Arrows}(\mathcal{D})
\mid H^i(f)\text{ is an isomorphism for all }i \in \mathbf{Z}\}
$$
Then $S$ is a saturated (see
Categories,
Definition \ref{categories-definition-saturated-multiplicative-system})
multiplicative system compatible with the
triangulated structure on $\mathcal{D}$.
\end{lemma}

\begin{proof}
We have to prove axioms MS1 -- MS6, see
Categories, Definitions \ref{categories-definition-multiplicative-system} and
\ref{categories-definition-saturated-multiplicative-system}
and
Definition \ref{definition-localization}.
MS1, MS4, and MS5 are direct from the definitions.
MS6 follows from TR3 and the long exact cohomology sequence
(\ref{equation-long-exact-cohomology-sequence}).
By
Lemma \ref{lemma-localization-conditions}
we conclude that MS2 holds. To finish the proof we have to show that
MS3 holds. To do this let $f, g : X \to Y$ be morphisms of $\mathcal{D}$,
and let $t : Z \to X$ be an element of $S$ such that $f \circ t = g \circ t$.
As $\mathcal{D}$ is additive this simply means that $a \circ t = 0$ with
$a = f - g$. Choose a distinguished triangle $(Z, X, Q, t, u, v)$ using
TR1 and TR2. Since $a \circ t = 0$ we see by
Lemma \ref{lemma-representable-homological}
there exists a morphism $i : Q \to Y$ such that $i \circ g = a$.
Finally, using TR1 again we can choose a triangle
$(Q, Y, W, i, j, k)$. Here is a picture
$$
\xymatrix{
Z \ar[r]_t & X \ar[r]_u \ar[d]^1 & Q \ar[r]_v \ar[d]^i & Z[1] \\
& X \ar[r]_a & Y \ar[d]^j \\
& & W
}
$$
OK, and now we apply the functors $H^i$ to this diagram.
Since $t \in S$ we see that $H^i(Q) = 0$ by the long exact cohomology
sequence (\ref{equation-long-exact-cohomology-sequence}).
Hence $H^i(j)$ is an isomorphism for all $i$ by the same argument,
i.e., $j \in S$. Finally, $j \circ a = j \circ i \circ u = 0$ as
$j \circ i = 0$. Thus $j \circ f = j \circ g$ and we see that LMS3 holds.
The proof of RMS3 is dual.
\end{proof}

\begin{proposition}
\label{proposition-construct-localization}
Let $\mathcal{D}$ be a pre-triangulated category. Let $S$ be a multiplicative
system compatible with the triangulated structure.
Then there exists a unique structure of a pre-triangulated category on
$S^{-1}\mathcal{D}$ such that $[1] \circ Q = Q \circ [1]$ and
the localization functor $Q : \mathcal{D} \to S^{-1}\mathcal{D}$ is exact.
Moreover, if $\mathcal{D}$ is a triangulated category, so is
$S^{-1}\mathcal{D}$.
\end{proposition}

\begin{proof}
We have seen that $S^{-1}\mathcal{D}$ is an additive category
and that the localization functor $Q$ is additive in
Homology, Lemma \ref{homology-lemma-localization-additive}.
It follows from MS5 that there is a unique additive auto-equivalence
$[1] : S^{-1}\mathcal{D} \to S^{-1}\mathcal{D}$
such that $Q \circ [1] = [1] \circ Q$ (equality of functors);
we omit the details. We say a triangle of $S^{-1}\mathcal{D}$
is distinguished if it is isomorphic to the image of a
distinguished triangle under the localization functor $Q$.

\medskip\noindent
Proof of TR1. The only thing to prove here is that if
$a : Q(X) \to Q(Y)$ is a morphism of $S^{-1}\mathcal{D}$, then
$a$ fits into a distinguished triangle. Write $a = Q(s)^{-1} \circ Q(f)$ for
some $s : Y \to Y'$ in $S$ and $f : X \to Y'$. Choose a distinguished
triangle $(X, Y', Z, f, g, h)$ in $\mathcal{D}$. Then we see that
$(Q(X), Q(Y), Q(Z), a, Q(g) \circ Q(s), Q(h))$ is a distinguished triangle
of $S^{-1}\mathcal{D}$.

\medskip\noindent
Proof of TR2. This is immediate from the definitions.

\medskip\noindent
Proof of TR3. Note that the existence of the dotted arrow which is
required to exist may be proven after replacing the two triangles
by isomorphic triangles. Hence we may assume given distinguished
triangles $(X, Y, Z, f, g, h)$ and $(X', Y', Z', f', g', h')$ of
$\mathcal{D}$ and a commutative diagram
$$
\xymatrix{
Q(X) \ar[r]_{Q(f)} \ar[d]_a & Q(Y) \ar[d]^b \\
Q(X') \ar[r]^{Q(f')} & Q(Y')
}
$$
in $S^{-1}\mathcal{D}$. Now we apply
Categories, Lemma \ref{categories-lemma-left-localization-diagram}
to find a morphism $f'' : X'' \to Y''$ in $\mathcal{D}$ and a commutative
diagram
$$
\xymatrix{
X \ar[d]_f \ar[r]_k & X'' \ar[d]^{f''} & X' \ar[d]^{f'} \ar[l]^s \\
Y \ar[r]^l & Y'' & Y' \ar[l]_t
}
$$
in $\mathcal{D}$ with $s, t \in S$ and $a = s^{-1}k$, $b = t^{-1}l$.
At this point we can use TR3 for $\mathcal{D}$ and MS6 to find
a commutative diagram
$$
\xymatrix{
X \ar[r] \ar[d]^k &
Y \ar[r] \ar[d]^l &
Z \ar[r] \ar[d]^m &
X[1] \ar[d]^{g[1]} \\
X'' \ar[r] &
Y'' \ar[r] &
Z'' \ar[r] &
X''[1] \\
X' \ar[r] \ar[u]_s &
Y' \ar[r] \ar[u]_t &
Z' \ar[r] \ar[u]_r &
X'[1] \ar[u]_{s[1]}
}
$$
with $r \in S$. It follows that setting $c = Q(r)^{-1}Q(m)$ we obtain
the desired morphism of triangles
$$
\xymatrix{
(Q(X), Q(Y), Q(Z), Q(f), Q(g), Q(h))
\ar[d]^{(a, b, c)} \\
(Q(X'), Q(Y'), Q(Z'), Q(f'), Q(g'), Q(h'))
}
$$

\medskip\noindent
This proves the first statement of the lemma. If $\mathcal{D}$ is also
a triangulated category, then we still have to prove TR4 in order to show
that $S^{-1}\mathcal{D}$ is triangulated as well. To do this we reduce by
Lemma \ref{lemma-easier-axiom-four}
to the following statement: Given composable morphisms
$a : Q(X) \to Q(Y)$ and $b : Q(Y) \to Q(Z)$ we have to produce
an octahedron after possibly replacing $Q(X), Q(Y), Q(Z)$ by isomorphic
objects. To do this we may first replace $Y$ by an object such that
$a = Q(f)$ for some morphism $f : X \to Y$ in $\mathcal{D}$. (More precisely,
write $a = s^{-1}f$ with $s : Y \to Y'$ in $S$ and $f : X \to Y'$. Then
replace $Y$ by $Y'$.) After this we similarly replace $Z$ by an object such
that $b = Q(g)$ for some morphism $g : Y \to Z$. Now we can find
distinguished triangles $(X, Y, Q_1, f, p_1, d_1)$,
$(X, Z, Q_2, g \circ f, p_2, d_2)$, and
$(Y, Z, Q_3, g, p_3, d_3)$ in $\mathcal{D}$ (by TR1), and
morphisms $a : Q_1 \to Q_2$ and $b : Q_2 \to Q_3$ as in TR4.
Then it is immediately verified that applying the functor $Q$ to
all these data gives a corresponding structure in $S^{-1}\mathcal{D}$.
\end{proof}

\noindent
The universal property of the localization of a triangulated category
is as follows (we formulate this for pre-triangulated categories, hence
it holds a fortiori for triangulated categories).

\begin{lemma}
\label{lemma-universal-property-localization}
Let $\mathcal{D}$ be a pre-triangulated category. Let $S$ be a multiplicative
system compatible with the triangulated structure. Let
$Q : \mathcal{D} \to S^{-1}\mathcal{D}$ be the localization functor, see
Proposition \ref{proposition-construct-localization}.
\begin{enumerate}
\item If $H : \mathcal{D} \to \mathcal{A}$ is a homological functor into
an abelian category $\mathcal{A}$ such that $H(s)$ is an isomorphism for
all $s \in S$, then the unique factorization
$H' : S^{-1}\mathcal{D} \to \mathcal{A}$ such that $H = H' \circ Q$ (see
Categories, Lemma \ref{categories-lemma-properties-left-localization})
is a homological functor too.
\item If $F : \mathcal{D} \to \mathcal{D}'$ is an exact functor into
a pre-triangulated category $\mathcal{D}'$ such that $F(s)$ is an isomorphism
for all $s \in S$, then the unique factorization
$F' : S^{-1}\mathcal{D} \to \mathcal{D}'$ such that $F = F' \circ Q$ (see
Categories, Lemma \ref{categories-lemma-properties-left-localization})
is an exact functor too.
\end{enumerate}
\end{lemma}

\begin{proof}
This lemma proves itself. Details omitted.
\end{proof}

\begin{lemma}
\label{lemma-localization-subcategory}
Let $\mathcal{D}$ be a pre-triangulated category and let
$\mathcal{D}' \subset \mathcal{D}$ be a full, pre-triangulated subcategory.
Let $S$ be a saturated multiplicative system of $\mathcal{D}$
compatible with the triangulated structure. Assume that for each $X$
in $\mathcal{D}$ there exists an $s : X' \to X$ in $S$ such that $X'$ is
an object of $\mathcal{D}'$. Then
$S' = S \cap \text{Arrows}(\mathcal{D}')$ is a saturated
multiplicative system compatible with the triangulated structure and
the functor
$$
(S')^{-1}\mathcal{D}' \longrightarrow S^{-1}\mathcal{D}
$$
is an equivalence of pre-triangulated categories.
\end{lemma}

\begin{proof}
Consider the quotient functor $Q : \mathcal{D} \to S^{-1}\mathcal{D}$
of Proposition \ref{proposition-construct-localization}.
Since $S$ is saturated we have that a morphism $f : X \to Y$ is in
$S$ if and only if $Q(f)$ is invertible, see
Categories, Lemma \ref{categories-lemma-what-gets-inverted}.
Thus $S'$ is the collection of arrows which are turned into isomorphisms
by the composition $\mathcal{D}' \to \mathcal{D} \to S^{-1}\mathcal{D}$.
Hence $S'$ is is a saturated multiplicative system compatible with the
triangulated structure by Lemma \ref{lemma-triangle-functor-localize}.
By Lemma \ref{lemma-universal-property-localization} we obtain the
exact functor $(S')^{-1}\mathcal{D}' \to S^{-1}\mathcal{D}$ of pre-triangulated
categories. By assumption this functor is essentially surjective.
Let $X', Y'$ be objects of $\mathcal{D}'$. By Categories, Remark
\ref{categories-remark-right-localization-morphisms-colimit} we have
$$
\Mor_{S^{-1}\mathcal{D}}(X', Y') =
\colim_{s : X \to X'\text{ in }S} \Mor_\mathcal{D}(X, Y')
$$
Our assumption implies that for any $s : X \to X'$ in $S$ we can
find a morphism $s' : X'' \to X$ in $S$ with $X''$ in $\mathcal{D}'$.
Then $s \circ s' : X'' \to X'$ is in $S'$. Hence the colimit above is
equal to
$$
\colim_{s'' : X'' \to X'\text{ in }S'} \Mor_{\mathcal{D}'}(X'', Y') =
\Mor_{(S')^{-1}\mathcal{D}'}(X', Y')
$$
This proves our functor is also fully faithful and the proof is complete.
\end{proof}

\noindent
The following lemma describes the kernel (see
Definition \ref{definition-kernel-category})
of the localization functor.

\begin{lemma}
\label{lemma-kernel-localization}
Let $\mathcal{D}$ be a pre-triangulated category. Let $S$ be a multiplicative
system compatible with the triangulated structure. Let $Z$ be an object
of $\mathcal{D}$. The following are equivalent
\begin{enumerate}
\item $Q(Z) = 0$ in $S^{-1}\mathcal{D}$,
\item there exists $Z' \in \Ob(\mathcal{D})$ such that
$0 : Z \to Z'$ is an element of $S$,
\item there exists $Z' \in \Ob(\mathcal{D})$ such that
$0 : Z' \to Z$ is an element of $S$, and
\item there exists an object $Z'$ and a distinguished triangle
$(X, Y, Z \oplus Z', f, g, h)$ such that $f \in S$.
\end{enumerate}
If $S$ is saturated, then these are also equivalent to
\begin{enumerate}
\item[(5)] the morphism $0 \to Z$ is an element of $S$,
\item[(6)] the morphism $Z \to 0$ is an element of $S$,
\item[(7)] there exists a distinguished triangle $(X, Y, Z, f, g, h)$
such that $f \in S$.
\end{enumerate}
\end{lemma}

\begin{proof}
The equivalence of (1), (2), and (3) is
Homology, Lemma \ref{homology-lemma-kernel-localization}.
If (2) holds, then $(Z'[-1], Z'[-1] \oplus Z, Z, (1, 0), (0, 1), 0)$
is a distinguished triangle (see
Lemma \ref{lemma-split})
with ``$0 \in S$''. By rotating we conclude that (4) holds.
If $(X, Y, Z \oplus Z', f, g, h)$ is a distinguished triangle with $f \in S$
then $Q(f)$ is an isomorphism hence $Q(Z \oplus Z') = 0$ hence $Q(Z) = 0$.
Thus (1) -- (4) are all equivalent.

\medskip\noindent
Next, assume that $S$ is saturated. Note that each of (5), (6), (7)
implies one of the equivalent conditions (1) -- (4). Suppose that
$Q(Z) = 0$. Then $0 \to Z$ is a morphism of $\mathcal{D}$ which becomes
an isomorphism in $S^{-1}\mathcal{D}$. According to
Categories, Lemma \ref{categories-lemma-what-gets-inverted}
the fact that $S$ is saturated implies that $0 \to Z$ is in $S$.
Hence (1) $\Rightarrow$ (5). Dually (1) $\Rightarrow$ (6).
Finally, if $0 \to Z$ is in $S$, then the triangle
$(0, Z, Z, 0, \text{id}_Z, 0)$ is distinguished by TR1 and TR2 and
is a triangle as in (7).
\end{proof}

\begin{lemma}
\label{lemma-limit-triangles}
Let $\mathcal{D}$ be a pre-triangulated category.
Let $S$ be a saturated multiplicative system in $\mathcal{D}$
that is compatible with the triangulated structure.
Let $(X, Y, Z, f, g, h)$ be a distinguished triangle in $\mathcal{D}$.
Consider the category of morphisms of triangles
$$
\mathcal{I} =
\{(s, s', s'') : (X, Y, Z, f, g, h) \to (X', Y', Z', f', g', h')
\mid s, s', s'' \in S\}
$$
Then $\mathcal{I}$ is a filtered category and the functors
$\mathcal{I} \to X/S$, $\mathcal{I} \to Y/S$, and $\mathcal{I} \to Z/S$
are cofinal.
\end{lemma}

\begin{proof}
We strongly suggest the reader skip the proof of this lemma and instead
work it out on a napkin.

\medskip\noindent
The first remark is that using rotation of distinguished triangles (TR2)
gives an equivalence of categories between $\mathcal{I}$ and the
corresponding category for the distinguished triangle
$(Y, Z, X[1], g, h, -f[1])$. Using this we see for example that if
we prove the functor $\mathcal{I} \to X/S$ is cofinal, then
the same thing is true for the functors $\mathcal{I} \to Y/S$ and
$\mathcal{I} \to Z/S$.

\medskip\noindent
Note that if $s : X \to X'$ is a morphism of $S$, then using
MS2 we can find $s' : Y \to Y'$ and $f' : X' \to Y'$ such that
$f' \circ s = s' \circ f$, whereupon we can use MS6 to complete
this into an object of $\mathcal{I}$. Hence the functor
$\mathcal{I} \to X/S$ is surjective on objects. Using rotation as above
this implies the same thing is true for the functors
$\mathcal{I} \to Y/S$ and $\mathcal{I} \to Z/S$.

\medskip\noindent
Suppose given objects $s_1 : X \to X_1 $ and $s_2 : X \to X_2$ in
$X/S$ and a morphism $a : X_1 \to X_2$ in $X/S$. Since $S$ is saturated,
we see that $a \in S$, see
Categories, Lemma \ref{categories-lemma-what-gets-inverted}.
By the argument of the previous paragraph we can complete
$s_1 : X \to X_1$ to an object
$(s_1, s'_1, s''_1) : (X, Y, Z, f, g, h) \to (X_1, Y_1, Z_1, f_1, g_1, h_1)$
in $\mathcal{I}$. Then we can repeat and find
$(a, b, c) : (X_1, Y_1, Z_1, f_1, g_1, h_1) \to (X_2, Y_2, Z_2, f_2, g_2, h_2)$
with $a, b, c \in S$ completing the given $a : X_1 \to X_2$.
But then $(a, b, c)$ is a morphism in $\mathcal{I}$.
In this way we conclude that the functor $\mathcal{I} \to X/S$ is
also surjective on arrows. Using rotation as above,
this implies the same thing is true for the functors
$\mathcal{I} \to Y/S$ and $\mathcal{I} \to Z/S$.

\medskip\noindent
The category $\mathcal{I}$ is nonempty as the identity provides an object.
This proves the condition (1) of the definition of a filtered category, see
Categories, Definition \ref{categories-definition-directed}.

\medskip\noindent
We check condition (2) of
Categories, Definition \ref{categories-definition-directed}
for the category $\mathcal{I}$. Suppose given objects
$(s_1, s'_1, s''_1) : (X, Y, Z, f, g, h) \to (X_1, Y_1, Z_1, f_1, g_1, h_1)$
and
$(s_2, s'_2, s''_2) : (X, Y, Z, f, g, h) \to (X_2, Y_2, Z_2, f_2, g_2, h_2)$
in $\mathcal{I}$. We want to find an object of $\mathcal{I}$
which is the target of an arrow from both
$(X_1, Y_1, Z_1, f_1, g_1, h_1)$ and $(X_2, Y_2, Z_2, f_2, g_2, h_2)$.
By Categories, Remark
\ref{categories-remark-left-localization-morphisms-colimit}
the categories $X/S$, $Y/S$, $Z/S$ are filtered.
Thus we can find $X \to X_3$ in $X/S$ and morphisms
$s : X_2 \to X_3$ and $a : X_1 \to X_3$. By the above we can find a morphism
$(s, s', s'') : (X_2, Y_2, Z_2, f_2, g_2, h_2) \to
(X_3, Y_3, Z_3, f_3, g_3, h_3)$ with $s', s'' \in S$.
After replacing $(X_2, Y_2, Z_2)$ by $(X_3, Y_3, Z_3)$ we may
assume that there exists a morphism $a : X_1 \to X_2$ in $X/S$.
Repeating the argument for $Y$ and $Z$ (by rotating as above)
we may assume there is a morphism
$a : X_1 \to X_2$ in $X/S$,
$b : Y_1 \to Y_2$ in $Y/S$, and
$c : Z_1 \to Z_2$ in $Z/S$.
However, these morphisms do not necessarily give rise to a morphism of
distinguished triangles. On the other hand, the necessary diagrams
do commute in $S^{-1}\mathcal{D}$. Hence we see (for example) that
there exists a morphism $s'_2 : Y_2 \to Y_3$ in $S$ such that
$s'_2 \circ f_2 \circ a = s'_2 \circ b \circ f_1$. Another replacement
of $(X_2, Y_2, Z_2)$ as above then gets us to the situation where
$f_2 \circ a = b \circ f_1$. Rotating and applying the same argument
two more times we see that we may assume $(a, b, c)$ is a morphism
of triangles. This proves condition (2).

\medskip\noindent
Next we check condition (3) of
Categories, Definition \ref{categories-definition-directed}.
Suppose $(s_1, s_1', s_1'') : (X, Y, Z) \to (X_1, Y_1, Z_1)$ and
$(s_2, s_2', s_2'') : (X, Y, Z) \to (X_2, Y_2, Z_2)$
are objects of $\mathcal{I}$, and suppose $(a, b, c), (a', b', c')$
are two morphisms between them. Since $a \circ s_1 = a' \circ s_1$
there exists a morphism $s_3 : X_2 \to X_3$ such that
$s_3 \circ a = s_3 \circ a'$. Using the surjectivity statement
we can complete this to a morphism of triangles
$(s_3, s_3', s_3'') : (X_2, Y_2, Z_2) \to (X_3, Y_3, Z_3)$
with $s_3, s_3', s_3'' \in S$. Thus
$(s_3 \circ s_2, s_3' \circ s_2', s_3'' \circ s_2'') :
(X, Y, Z) \to (X_3, Y_3, Z_3)$ is also an object of $\mathcal{I}$
and after composing the maps $(a, b, c), (a', b', c')$ with
$(s_3, s_3', s_3'')$ we obtain $a = a'$. By rotating we may do the
same to get $b = b'$ and $c = c'$.

\medskip\noindent
Finally, we check that $\mathcal{I} \to X/S$ is cofinal, see
Categories, Definition \ref{categories-definition-cofinal}.
The first condition is true as the functor is surjective.
Suppose that we have an object $s : X \to X'$ in $X/S$ and
two objects
$(s_1, s'_1, s''_1) : (X, Y, Z, f, g, h) \to (X_1, Y_1, Z_1, f_1, g_1, h_1)$
and
$(s_2, s'_2, s''_2) : (X, Y, Z, f, g, h) \to (X_2, Y_2, Z_2, f_2, g_2, h_2)$
in $\mathcal{I}$ as well as morphisms $t_1 : X' \to X_1$ and
$t_2 : X' \to X_2$ in $X/S$. By property (2) of $\mathcal{I}$
proved above we can find morphisms
$(s_3, s'_3, s''_3) : (X_1, Y_1, Z_1, f_1, g_1, h_1) \to
(X_3, Y_3, Z_3, f_3, g_3, h_3)$
and
$(s_4, s'_4, s''_4) : (X_2, Y_2, Z_2, f_2, g_2, h_2) \to
(X_3, Y_3, Z_3, f_3, g_3, h_3)$ in $\mathcal{I}$.
We would be done if the compositions
$X' \to X_1 \to X_3$ and $X' \to X_2 \to X_3$ were equal
(see displayed equation in 
Categories, Definition \ref{categories-definition-cofinal}).
If not, then, because $X/S$ is filtered, we can choose
a morphism $X_3 \to X_4$ in $X/S$ such that the compositions
$X' \to X_1 \to X_3 \to X_4$ and $X' \to X_2 \to X_3 \to X_4$ are equal.
Then we finally complete $X_3 \to X_4$ to a morphism
$(X_3, Y_3, Z_3) \to (X_4, Y_4, Z_4)$ in $\mathcal{I}$
and compose with that morphism to see that the result is true.
\end{proof}









